% Public source snapshot of the instructor's Module 11 lecture notes. Executable
% implementations, diagram source, and external images are omitted under the
% boundary in sources/README.md. Headings and claims are preserved for provenance,
% not endorsed as reproduced results; the book re-derives equations and checks
% numerical claims against primary papers before use.
\documentclass{article}
\usepackage{graphicx}
\usepackage{amsmath}
\usepackage{amssymb}
\usepackage{listings}
\usepackage[margin=1in]{geometry}
\usepackage{xcolor}
\usepackage{tikz}
\usetikzlibrary{shapes, arrows, positioning, calc}
\usepackage{tcolorbox}
\usepackage{booktabs}

% Define colors
\definecolor{codegreen}{rgb}{0,0.6,0}
\definecolor{codegray}{rgb}{0.5,0.5,0.5}
\definecolor{codepurple}{rgb}{0.58,0,0.82}
\definecolor{backcolour}{rgb}{0.95,0.95,0.92}

% Python style
\lstdefinestyle{pythonstyle}{
    backgroundcolor=\color{backcolour},
    commentstyle=\color{codegreen},
    keywordstyle=\color{magenta},
    numberstyle=\tiny\color{codegray},
    stringstyle=\color{codepurple},
    basicstyle=\ttfamily\footnotesize,
    breakatwhitespace=false,
    breaklines=true,
    captionpos=b,
    keepspaces=true,
    numbers=left,
    numbersep=5pt,
    showspaces=false,
    showstringspaces=false,
    showtabs=false,
    tabsize=2
}
\lstset{style=pythonstyle}

\title{Lecture 11.1: Adapting Large Language Models Without Fine-Tuning\\
In-Context Learning, Prompt Engineering, and RAG}
\author{Deep Learning Class}
\date{\today}

\begin{document}
\maketitle

\begin{tcolorbox}[colback=gray!5!white, colframe=gray!60!black, title=Public Snapshot Boundary]
This provenance snapshot preserves the instructor-authored lecture spine. Executable
implementations, diagram source, and external images have been omitted; their original
locations remain marked in the source. Product and numerical claims are historical
lecture notes and are independently checked before publication in the book.
\end{tcolorbox}

\begin{tcolorbox}[colback=blue!5!white, colframe=blue!75!black, title=Lecture Overview]
Today's focus: How to adapt LLMs to specific tasks \textit{without} modifying their parameters.

\textbf{Central thesis}: When you can get away with no fine-tuning, you absolutely should!

\textbf{Topics}:
\begin{enumerate}
\item The adaptation challenge: generality vs specificity
\item In-Context Learning (ICL): Learning without parameter updates
\item Prompt engineering: The art and science of effective prompts
\item Retrieval-Augmented Generation (RAG): Grounding LLMs in external knowledge
\item Why hesitate to fine-tune? Costs, risks, and hidden dangers
\item Decision framework: When to use each approach
\item Practical examples with Llama 3.1-8B
\end{enumerate}
\end{tcolorbox}

\section{The Adaptation Challenge}

\subsection{The Generality-Specificity Tension}

Large Language Models like GPT-4, Claude, and Llama 3.1 represent a remarkable achievement: models with general-purpose natural language understanding and generation. This proficiency comes from pre-training on trillions of tokens of text.

\textbf{The challenge}: This very generality presents a fundamental problem for practical applications.

\begin{center}
% [Diagram source omitted from the public snapshot; the book uses independently drawn figures.]
\end{center}

\textbf{Key limitations of base models}:
\begin{itemize}
\item \textbf{Static knowledge}: Fixed training cutoff date (e.g., Llama 3 trained on data through 2023)
\item \textbf{Diffuse knowledge}: Broad but not deep on specialized domains
\item \textbf{No private data}: Cannot access proprietary company information
\item \textbf{Lack of specificity}: May not follow exact style, format, or behavioral requirements
\end{itemize}

\subsection{Two Paradigms for Adaptation}

How do we bridge this gap? There are two fundamentally different approaches:

\begin{center}
% [Diagram source omitted from the public snapshot; the book uses independently drawn figures.]
\end{center}

\textbf{Context Augmentation}: Treat the LLM as a fixed, powerful reasoning engine. Guide its output by providing information, examples, and instructions in the prompt at inference time. The model's weights never change.

\textbf{Parameter Modification}: Adjust the internal weights of the model through additional training on task-specific data, embedding new knowledge or behaviors directly into the model.

\begin{tcolorbox}[colback=yellow!5!white, colframe=orange!50!black, title=Central Principle]
\textbf{``When you can get away with no fine-tuning, you absolutely should!''}

This lecture will explain why this principle holds and show you powerful techniques that often eliminate the need for fine-tuning entirely.
\end{tcolorbox}

\section{In-Context Learning: The Remarkable Capability}

\subsection{What is In-Context Learning?}

\textbf{Definition}: In-Context Learning (ICL) is the ability of large language models to learn a new task from examples or instructions provided in the prompt, \textit{without any parameter updates}.

\textbf{Key characteristics}:
\begin{itemize}
\item Model weights remain \textbf{completely frozen}
\item ``Learning'' is transient, lasting only for that inference
\item The prompt itself serves as the ``training'' for the task
\item No gradient descent, no backpropagation, no optimization
\end{itemize}

\begin{center}
% [Diagram source omitted from the public snapshot; the book uses independently drawn figures.]
\end{center}

\subsection{How ICL Works: The Mechanism}

\textbf{Core mechanism}: Pattern recognition and analogical reasoning through the attention mechanism.

When you provide examples in the prompt, the model's multi-head attention analyzes these demonstrations to:
\begin{enumerate}
\item Identify the underlying task structure
\item Recognize input-output patterns
\item Infer the desired format and style
\item Extrapolate to new, unseen queries
\end{enumerate}

\textbf{Example}: Sentiment classification with ICL

\begin{tcolorbox}[colback=gray!5!white, colframe=gray!50!black, title=ICL Prompt]
\textbf{Task}: Classify the sentiment of movie reviews as positive or negative.

\textbf{Examples}:
\begin{itemize}
\item Review: ``This movie was absolutely fantastic! I loved every minute.''\\
Sentiment: Positive
\item Review: ``Terrible film. Waste of time and money.''\\
Sentiment: Negative
\item Review: ``An unexpected masterpiece with brilliant acting.''\\
Sentiment: Positive
\end{itemize}

\textbf{New query}:\\
Review: ``The plot was confusing and the pacing was too slow.''\\
Sentiment: ?
\end{tcolorbox}

The model recognizes the pattern and responds: ``Negative''

\textbf{What makes this remarkable}:
\begin{itemize}
\item No training data collection or annotation
\item No model training or optimization
\item Instant adaptation at inference time
\item Same model can handle countless tasks by changing the prompt
\end{itemize}

\subsection{The Scaling Law for ICL}

\begin{tcolorbox}[colback=blue!5!white, colframe=blue!75!black, title=Critical Finding: ICL and Model Scale]
The effectiveness of In-Context Learning is \textbf{strongly dependent on model scale}.

\textbf{Empirical observation}: Larger models, pre-trained on more data, exhibit dramatically better ICL performance.

\textbf{Why?} Extensive pre-training endows large models with sophisticated understanding of language patterns, enabling them to recognize and utilize contextual cues more effectively.
\end{tcolorbox}

\textbf{Implications for practice}:
\begin{itemize}
\item Small models ($<$1B parameters): ICL capabilities are limited
\item Medium models (7-13B parameters like Llama 3.1-8B): Strong ICL for many tasks
\item Large models ($>$70B parameters): Exceptional ICL, approaching fine-tuned performance
\end{itemize}

\textbf{Good news}: Modern open-weight models like Llama 3.1-8B have reached the scale where ICL is highly effective for a wide range of applications.

\section{Prompt Engineering: Crafting Effective Prompts}

\subsection{The Importance of Prompt Design}

\textbf{Key insight}: The model's output \textit{heavily depends} on how the input is phrased and structured.

Small changes in wording, example ordering, or formatting can significantly affect results. This makes prompt engineering both an art and a science.

\begin{center}
% [Diagram source omitted from the public snapshot; the book uses independently drawn figures.]
\end{center}

\textbf{Prompt engineering}: The practice of systematically designing effective prompts to guide the model's behavior and maximize task performance.

\subsection{Core Prompting Strategies}

\subsubsection{1. Zero-Shot Prompting}

\textbf{Definition}: Give the model an instruction or question with no examples---just a description of the task.

\textbf{Example}:
% [Implementation omitted from the public snapshot: independent provenance was not established. See sources/README.md.]

\textbf{When to use}: For simple, well-defined tasks where the model likely has seen many similar examples during pre-training.

\subsubsection{2. Few-Shot Prompting}

\textbf{Definition}: Provide one or more demonstration examples of input-output pairs so the model can mimic the pattern.

\textbf{Example}:
% [Implementation omitted from the public snapshot: independent provenance was not established. See sources/README.md.]

\textbf{Best practices}:
\begin{itemize}
\item Use 2-5 diverse, high-quality examples
\item Ensure examples cover edge cases
\item Maintain consistent formatting across examples
\item Order matters: sometimes quality examples should go first
\end{itemize}

\subsubsection{3. Chain-of-Thought (CoT) Prompting}

\textbf{Breakthrough technique}: Dramatically improves performance on complex reasoning tasks by including intermediate reasoning steps in the examples.

\textbf{Key insight}: Standard prompting asks for a direct answer, forcing the model to compute the solution in one step. CoT prompting encourages the model to ``think step by step'', decomposing the problem.

\begin{tcolorbox}[colback=yellow!5!white, colframe=orange!50!black, title=Power of Chain-of-Thought]
\textbf{Impact}: CoT prompting can improve accuracy by up to +18\% on arithmetic reasoning tasks and even more on complex multi-step problems.

\textbf{Why it works}:
\begin{itemize}
\item Allocates more computation to each reasoning step
\item Makes the reasoning process explicit and verifiable
\item Helps the model avoid shortcuts that lead to errors
\item Provides interpretability into the model's thought process
\end{itemize}
\end{tcolorbox}

\textbf{Example: Standard vs CoT}

\textbf{Standard prompting}:
\begin{tcolorbox}[colback=gray!5!white, colframe=gray!50!black]
Q: Roger has 5 tennis balls. He buys 2 more cans of tennis balls. Each can has 3 tennis balls. How many tennis balls does he have now?

A: 11
\end{tcolorbox}

\textbf{Chain-of-Thought prompting}:
\begin{tcolorbox}[colback=green!5!white, colframe=green!50!black]
Q: Roger has 5 tennis balls. He buys 2 more cans of tennis balls. Each can has 3 tennis balls. How many tennis balls does he have now?

A: Roger started with 5 balls. 2 cans of 3 tennis balls each is $2 \times 3 = 6$ tennis balls. $5 + 6 = 11$. The answer is 11.
\end{tcolorbox}

\textbf{Advanced CoT Variants}:

\begin{itemize}
\item \textbf{Zero-Shot CoT}: Simply append ``Let's think step by step'' to the query---no examples needed! Surprisingly effective for large models.

\item \textbf{Self-Consistency}: Generate multiple reasoning chains for the same question (using sampling), then take a majority vote among the final answers. More robust to occasional reasoning errors.

\item \textbf{Auto-CoT}: Automatically generate reasoning chain examples by clustering questions and generating chains for representative samples.
\end{itemize}

\subsection{Automated Prompt Optimization}

\textbf{Challenge}: Manually crafting optimal prompts can be time-consuming and requires expertise.

\textbf{Solution}: Use learning algorithms or auxiliary models to automatically improve prompts.

\begin{center}
% [Diagram source omitted from the public snapshot; the book uses independently drawn figures.]
\end{center}

\textbf{Key approaches}:

\begin{enumerate}
\item \textbf{AutoPrompt}: Search in the space of possible prompts (even at token level) using gradient signals. Can discover prompts that match fine-tuned model performance without parameter updates.

\item \textbf{RL-based optimization}: Train a small neural policy network to generate prompts that maximize the target LLM's performance reward.

\item \textbf{LLM-as-optimizer}: Use a separate LLM to iteratively refine prompts based on feedback.
\end{enumerate}

\textbf{Key insight}: We can often achieve better task performance by improving the prompt itself---using machine learning or programmatic search---rather than adjusting the large model's weights.

\subsection{Practical Prompt Engineering with Llama 3.1-8B}

% [Implementation omitted from the public snapshot: independent provenance was not established. See sources/README.md.]

% [Implementation omitted from the public snapshot: independent provenance was not established. See sources/README.md.]

\section{Retrieval-Augmented Generation (RAG)}

\subsection{The Knowledge Problem}

\textbf{Limitation}: Even with excellent ICL and prompting, the model can only work with:
\begin{itemize}
\item Knowledge from its pre-training (cutoff date: 2023 for Llama 3)
\item Information you can fit in the context window
\end{itemize}

\textbf{Problem scenarios}:
\begin{itemize}
\item Need current information (news, stock prices, recent events)
\item Large proprietary knowledge base (company documents, medical records)
\item Domain-specific information not well-represented in pre-training
\item Verifiable, attributed answers (citations required)
\end{itemize}

\textbf{Solution}: Retrieval-Augmented Generation (RAG)

\subsection{RAG Architecture and Pipeline}

\textbf{Core principle}: Combine the LLM's reasoning abilities with an external, updatable knowledge base. Ground the model's responses in retrieved, authoritative information.

\begin{center}
% [Diagram source omitted from the public snapshot; the book uses independently drawn figures.]
\end{center}

\subsection{RAG Pipeline Steps}

\subsubsection{Step 1: Indexing and Data Preparation (Offline)}

\textbf{Input}: External corpus (PDFs, documents, websites, databases)

\textbf{Process}:
\begin{enumerate}
\item \textbf{Parse}: Extract text from various formats
\item \textbf{Chunk}: Segment into smaller pieces
\begin{itemize}
\item Too large: Too general, may overflow context window
\item Too small: Loses semantic context
\item Typical: 256-512 tokens per chunk with overlap
\end{itemize}
\item \textbf{Embed}: Transform each chunk into a high-dimensional vector using an embedding model (e.g., sentence-transformers)
\item \textbf{Store}: Save vectors in a vector database (e.g., FAISS, Pinecone, Weaviate, Chroma)
\end{enumerate}

\textbf{Key property}: In the vector space, semantically similar chunks are located close to each other, enabling semantic search rather than just keyword matching.

\subsubsection{Step 2: Retrieval at Query Time (Online)}

\textbf{Input}: User query

\textbf{Process}:
\begin{enumerate}
\item \textbf{Embed query}: Convert query to vector using the \textit{same} embedding model
\item \textbf{Similarity search}: Find top-$k$ most similar chunks using cosine similarity or dot product
\item \textbf{(Optional) Re-rank}: Use a re-ranking model to further refine relevance
\end{enumerate}

\textbf{Output}: Top-$k$ most relevant text chunks from the knowledge base

\subsubsection{Step 3: Prompt Augmentation and Generation (Online)}

\textbf{Input}: Original query + retrieved chunks

\textbf{Process}:
\begin{enumerate}
\item \textbf{Synthesize prompt}: Combine query and retrieved content
\item \textbf{Generate}: Pass augmented prompt to LLM (e.g., Llama 3.1-8B)
\item \textbf{Output}: Comprehensive, grounded answer with optional source attribution
\end{enumerate}

\subsection{The Core Value Proposition of RAG}

\begin{center}
\begin{tabular}{lll}
\toprule
\textbf{Challenge} & \textbf{Fine-Tuning Approach} & \textbf{RAG Approach} \\
\midrule
Hallucinations & Reduces but doesn't eliminate & Grounds in verifiable sources \\
Data freshness & Requires retraining & Update knowledge base instantly \\
Transparency & Black box & Can cite source documents \\
Privacy/Security & Embeds in weights & Data stays in secure database \\
Cost & High (retraining) & Moderate (maintain database) \\
\bottomrule
\end{tabular}
\end{center}

\subsubsection{1. Mitigation of Hallucinations}

\textbf{Problem}: LLMs can generate factually incorrect or nonsensical information (``hallucination'').

\textbf{RAG solution}: By grounding every response in specific retrieved information, RAG significantly reduces hallucination. The model is instructed to base its answer on provided context.

\subsubsection{2. Ensuring Data Freshness}

\textbf{Problem}: Foundation models have fixed knowledge cutoff dates.

\textbf{RAG solution}: The external knowledge base can be continuously updated. No need to retrain the model---just update the vector database.

\textbf{Example}: Llama 3 has a 2023 cutoff. Using RAG, it can answer questions about 2024 or 2025 events by retrieving recent documents.

\subsubsection{3. Source Attribution and Transparency}

\textbf{RAG advantage}: Can present retrieved source documents as citations, allowing users to verify information.

\textbf{Fine-tuning limitation}: Knowledge is opaquely embedded in weights---impossible to trace the origin of specific information.

\subsubsection{4. Enhanced Security and Privacy}

\textbf{For enterprises with proprietary data}:
\begin{itemize}
\item Store confidential data in secure, on-premise vector database
\item LLM only sees small, relevant snippets at query time
\item Data never used to train or update model parameters
\item Prevents sensitive information from leaking to third-party models
\end{itemize}

\subsection{RAG and Large Context Windows: A Game Changer}

\begin{tcolorbox}[colback=blue!5!white, colframe=blue!75!black, title=The 128K Context Window Revolution]
\textbf{Llama 3.1}: 128,000-token context window (equivalent to $\sim$300 pages of text)

\textbf{Impact}: Fundamentally reshapes the RAG vs fine-tuning calculus.

\textbf{Historical constraint}: Early models had 2K-4K token windows, making it impossible to provide sufficient context in the prompt. Fine-tuning was often the only option for knowledge injection.

\textbf{New reality}: With 128K tokens, we can:
\begin{itemize}
\item Include hundreds of retrieved documents in a single prompt
\item Use ``context stuffing''---provide entire knowledge base directly
\item Eliminate the need for fine-tuning for many knowledge tasks
\end{itemize}
\end{tcolorbox}

\textbf{Strategic implications}:
\begin{itemize}
\item For \textbf{large, dynamic corpora}: Use RAG with retrieval
\item For \textbf{small, static document sets}: Direct context injection (no RAG needed)
\item \textbf{Reserve fine-tuning} for behavior/style modification, not knowledge transfer
\end{itemize}

\subsection{Practical RAG Implementation with Llama 3.1-8B}

% [Implementation omitted from the public snapshot: independent provenance was not established. See sources/README.md.]

% [Implementation omitted from the public snapshot: independent provenance was not established. See sources/README.md.]

\section{Why Hesitate to Fine-Tune?}

We've seen powerful techniques for adapting LLMs without touching their parameters. But why should we prefer these methods over fine-tuning? Let's examine the hidden costs and dangers of parameter modification.

\subsection{Reason 1: Catastrophic Forgetting}

\textbf{Definition}: When fine-tuned on a new task, the model often loses proficiency on previously learned tasks. The fine-tuning process overwrites knowledge acquired during pre-training.

\begin{center}
% [Diagram source omitted from the public snapshot; the book uses independently drawn figures.]
\end{center}

\textbf{Example}: Fine-tune Llama 3.1-8B extensively on legal documents:
\begin{itemize}
\item[$+$] Improves legal reasoning
\item[$-$] May lose ability to write creative stories
\item[$-$] May degrade at code generation
\item[$-$] May lose multilingual capabilities
\end{itemize}

\textbf{Mechanism}: During fine-tuning, gradient updates shift weights to optimize for the new task. This can move weights away from optimal points for pre-trained tasks, effectively ``overwriting'' that knowledge.

\textbf{Implications}:
\begin{itemize}
\item Fine-tuning is a \textbf{trade-off}: specialized expertise at the cost of generalist capability
\item May need to maintain multiple model versions (base + fine-tuned) for different use cases
\item ICL and RAG preserve the original model's versatility
\end{itemize}

\subsection{Reason 2: The Alignment Tax}

\begin{tcolorbox}[colback=red!5!white, colframe=red!75!black, title=The Alignment Tax]
\textbf{Definition}: The process of aligning an LLM with human preferences (via RLHF, DPO, etc.) can lead to degradation in the model's performance on standard benchmarks.

\textbf{Trade-off}: As a model becomes more ``aligned'' (helpful, harmless, honest), its raw capabilities may decline.
\end{tcolorbox}

\textbf{Mechanism}:
\begin{itemize}
\item Alignment constrains the model's output space
\item Teaching caution and safety may reduce creative reasoning
\item Model may overfit to biases in human preference data
\item Knowledge forgetting occurs during alignment training
\end{itemize}

\textbf{Implication}: Every step to make a model more polished through fine-tuning may simultaneously make it less capable in fundamental ways.

\subsection{Reason 3: Compromised Safety Guardrails}

\begin{tcolorbox}[colback=red!10!white, colframe=red!75!black, title=Critical Security Risk]
\textbf{Alarming finding}: Fine-tuning can easily and systematically dismantle safety guardrails that prevent harmful outputs.

\textbf{Research from Stanford, Princeton, and IBM}:
\begin{itemize}
\item Fine-tuning on as few as \textbf{10 harmful examples} can ``jailbreak'' models like GPT-3.5 or Llama 2
\item \textbf{Cost}: Only \$0.20 on a commercial API
\item Even fine-tuning on \textbf{benign datasets} can unintentionally degrade safety
\end{itemize}

\textbf{Mechanism}: Fine-tuning doesn't teach new harmful behaviors---it \textbf{removes} safety guardrails, revealing latent dangerous capabilities.
\end{tcolorbox}

\textbf{Example}: Llama 3.1-8B-Instruct has been carefully aligned to refuse:
\begin{itemize}
\item Instructions for building weapons
\item Generating malware
\item Creating hate speech
\item Harmful medical advice
\end{itemize}

\textbf{Risk}: Fine-tuning for any purpose may weaken these protections.

\textbf{Contrast with RAG/ICL}: These methods do not modify parameters, so they \textbf{do not interfere} with core safety alignment.

\subsection{Reason 4: Expertise and Scale Required}

\textbf{Challenge 1: Computational Cost}
\begin{itemize}
\item Full fine-tuning of 8B model: Requires high-end GPU (A100/H100)
\item Training time: Hours to days depending on dataset size
\item Cost: \$100s to \$1000s per training run
\end{itemize}

\textbf{Challenge 2: Data Requirements}
\begin{itemize}
\item Need high-quality, labeled dataset
\item Typical: 1000s to 100,000s of examples
\item Manual annotation by domain experts: \$100,000+ for complex domains
\end{itemize}

\textbf{Challenge 3: Technical Expertise}
\begin{itemize}
\item Requires deep ML/NLP expertise
\item Hyperparameter tuning (learning rate, batch size, epochs)
\item Evaluation and debugging
\item Risk of overfitting or misalignment if done incorrectly
\end{itemize}

\textbf{Contrast}: Prompt engineering and RAG require more software engineering and data management skills than deep learning expertise---lower barrier to entry.

\subsection{Reason 5: Maintenance and Update Complexity}

\textbf{Fine-tuning creates a fork}:
\begin{itemize}
\item Creates a new model instance specific to your task
\item When base model is updated (bug fixes, improvements), your fine-tuned version doesn't automatically benefit
\item Must re-fine-tune on new base model to incorporate improvements
\end{itemize}

\textbf{RAG/Prompting advantage}:
\begin{itemize}
\item Can upgrade to newer/better base model effortlessly
\item Just prompt the new model the same way or plug retrieval into it
\item No retraining required
\end{itemize}

\subsection{Reason 6: Models Are Optimized by Labs with Expertise}

\begin{tcolorbox}[colback=yellow!5!white, colframe=orange!50!black, title=Trust the Experts]
State-of-the-art LLMs like Llama 3.1 are the result of:
\begin{itemize}
\item Years of research and engineering
\item Billions of dollars in compute
\item Careful optimization by world-class ML teams
\item Extensive safety alignment and red-teaming
\end{itemize}

\textbf{Risk}: By modifying weights without the same level of rigor, you may:
\begin{itemize}
\item Deteriorate overall performance
\item Violate safety alignment
\item Introduce unexpected behaviors
\end{itemize}

\textbf{Principle}: If you're not Meta/OpenAI/Anthropic with massive resources, be cautious about changing what they've built.
\end{tcolorbox}

\section{Decision Framework: When to Use Each Approach}

\subsection{The Stepwise Approach}

Follow this progression---\textbf{always try less invasive methods first}:

\begin{center}
% [Diagram source omitted from the public snapshot; the book uses independently drawn figures.]
\end{center}

\subsection{Comparative Analysis}

\begin{center}
\small
\begin{tabular}{p{2.5cm}p{3.5cm}p{3.5cm}p{3.5cm}}
\toprule
\textbf{Dimension} & \textbf{Prompting/ICL} & \textbf{RAG} & \textbf{Fine-Tuning} \\
\midrule
\textbf{Goal} & Task guidance, reasoning & Knowledge injection & Behavior/style change \\
\midrule
\textbf{Data freshness} & N/A & Excellent (real-time) & Poor (static) \\
\midrule
\textbf{Hallucination} & Moderate & Low (grounded) & Moderate-Low \\
\midrule
\textbf{Transparency} & High & High (citations) & Low (black box) \\
\midrule
\textbf{Safety risk} & None & None & High \\
\midrule
\textbf{Forgetting} & None & None & High \\
\midrule
\textbf{Cost} & Very low & Low-Moderate & High \\
\midrule
\textbf{Expertise} & Low & Moderate & High \\
\midrule
\textbf{Latency} & Low & Moderate & Low \\
\midrule
\textbf{Maintenance} & Easy & Easy (update DB) & Hard (retrain) \\
\bottomrule
\end{tabular}
\end{center}

\subsection{When Fine-Tuning IS Appropriate}

Fine-tuning is the right choice when:

\begin{enumerate}
\item \textbf{Behavioral modification}: Need to change how the model acts, not what it knows
\begin{itemize}
\item Example: Always output in JSON format
\item Example: Adopt specific brand personality/tone
\item Example: Follow domain-specific reasoning patterns
\end{itemize}

\item \textbf{Style enforcement}: Need consistent style that prompting can't reliably achieve
\begin{itemize}
\item Example: Medical report writing style
\item Example: Legal document formatting
\end{itemize}

\item \textbf{Skill acquisition}: Teaching complex, specialized skills
\begin{itemize}
\item Example: Proprietary programming language
\item Example: Domain-specific reasoning (e.g., radiology interpretation)
\end{itemize}

\item \textbf{Prompt compression}: When optimal prompt is too long for context window
\begin{itemize}
\item Can ``distill'' long instructions into model weights
\end{itemize}
\end{enumerate}

\textbf{Important}: Even when fine-tuning is appropriate, use \textbf{Parameter-Efficient Fine-Tuning (PEFT)} methods like LoRA (topic of next lecture) to minimize risks.

\section{Summary and Looking Ahead}

\subsection{Key Takeaways}

\begin{enumerate}
\item \textbf{In-Context Learning} is a remarkable capability that allows LLMs to adapt to new tasks without parameter updates. Modern models like Llama 3.1-8B have strong ICL abilities.

\item \textbf{Prompt engineering} is the first line of attack for any adaptation task. Simple techniques like few-shot prompting and Chain-of-Thought can achieve impressive results.

\item \textbf{RAG} is the superior method for knowledge injection. It provides:
\begin{itemize}
\item Fresh, updatable information
\item Reduced hallucination
\item Source attribution and transparency
\item Better security for proprietary data
\end{itemize}

\item \textbf{Large context windows} (128K tokens in Llama 3.1) fundamentally change the game, making fine-tuning for knowledge transfer largely obsolete.

\item \textbf{Fine-tuning carries serious risks}:
\begin{itemize}
\item Catastrophic forgetting
\item Alignment tax
\item Safety guardrail degradation
\item High cost and expertise requirements
\item Maintenance complexity
\end{itemize}

\item \textbf{General principle}: ``When you can get away with no fine-tuning, you absolutely should!''
\end{enumerate}

\subsection{The Adaptation Hierarchy}

\begin{center}
% [Diagram source omitted from the public snapshot; the book uses independently drawn figures.]
\end{center}

\subsection{Next Lecture: When We Must Fine-Tune}

In the next lecture, we'll cover Parameter-Efficient Fine-Tuning (PEFT):

\begin{itemize}
\item When fine-tuning is truly necessary
\item The baseline: Full Fine-Tuning and its costs
\item PEFT taxonomy: Additive, Selective, and Prompt-Based methods
\item LoRA: Low-Rank Adaptation in depth
\item Practical implementation with Llama 3.1-8B
\item Hybrid approaches: Combining RAG with fine-tuning
\end{itemize}

\begin{tcolorbox}[colback=blue!5!white, colframe=blue!75!black, title=The Complete Picture]
\textbf{Best practice}: Use a modular, hybrid approach
\begin{itemize}
\item Start with prompting and RAG (today's lecture)
\item If necessary, use minimal PEFT (next lecture)
\item Combine approaches for optimal results
\end{itemize}

The future of LLM applications is not about choosing between fine-tuning and prompting---it's about intelligently combining them in a sophisticated system architecture.
\end{tcolorbox}

\section{Further Reading}

\subsection{In-Context Learning}
\begin{itemize}
\item Brown et al. (2020). ``Language Models are Few-Shot Learners.'' (GPT-3 paper)
\item IBM Research. ``What is In-Context Learning?''
\item Wei et al. (2022). ``Emergent Abilities of Large Language Models.''
\end{itemize}

\subsection{Prompt Engineering}
\begin{itemize}
\item Wei et al. (2022). ``Chain-of-Thought Prompting Elicits Reasoning in Large Language Models.''
\item Kojima et al. (2022). ``Large Language Models are Zero-Shot Reasoners.''
\item Shin et al. (2020). ``AutoPrompt: Eliciting Knowledge from Language Models with Automatically Generated Prompts.''
\end{itemize}

\subsection{Retrieval-Augmented Generation}
\begin{itemize}
\item Lewis et al. (2020). ``Retrieval-Augmented Generation for Knowledge-Intensive NLP Tasks.''
\item Gao et al. (2023). ``Retrieval-Augmented Generation for Large Language Models: A Survey.''
\end{itemize}

\subsection{Fine-Tuning Risks}
\begin{itemize}
\item Qi et al. (2023). ``Fine-tuning Aligned Language Models Compromises Safety, Even When Users Do Not Intend To!'' (Stanford/Princeton/IBM)
\item Kirkpatrick et al. (2017). ``Overcoming Catastrophic Forgetting in Neural Networks.''
\end{itemize}

\end{document}
